% Owner: WP4 (Ben) with WP2 (Angelo). Tasks: A-14, R-16.
% Required: end-to-end component diagram separating client UI, FastAPI orchestration,
% inference engine, vector store and relational database; request lifecycle of a
% multi-hop query with token transfer and guardrail validation steps.
\section{System Architecture and Component Decoupling}
\label{sec:architecture}

\paragraph{Architecture.}
\todo{Diagram (A-14): Streamlit UI, FastAPI, Ollama on Modal, ParadeDB with pgvector and BM25, relational tables. Put the file in figures/.}

\paragraph{Execution flow.}
For an S3 question about the most cited authors on a topic, FastAPI first runs
\texttt{check\_input}: Prompt Guard, the topic/write-intent gate and PII masking.
A refusal ends the request before planning or database retrieval.
The model produces a validated scenario plan; explicit year and result-limit filters
are reconciled in code. Filtered dense and BM25 retrieval selects topic papers and
their passages. The evidence guard removes flagged sources before SQL and generation.
SQL follows outgoing citation edges from the surviving paper seeds to cited papers
and their authors. Each distinct citing/cited-paper pair gives every cited-paper
coauthor one credit; these are subset counts, not global citation totals.
The resulting table, its parameters, dataset revision and query ID are emitted directly
by the application, before generated sentences. For a plain ranking request the
application can finish using the table and a cited counting-rule sentence. For other
S3 requests the model summarizes paper evidence and does not calculate the ranking.

The OpenAI-compatible model stream carries text deltas to a bounded server-side
sentence buffer. Raw deltas are not forwarded to the UI. Each complete sentence,
including the final punctuated tail, passes citation-structure checks and
\texttt{check\_sentence}: numerical consistency against cited evidence, then NLI for
non-SQL sources. SQL evidence receives the numerical check. Accepted sentences are
PII-masked, rendered with citation objects and released as sentence SSE events;
rejected sentences produce a removal status. Timeout or provider failure produces
a terminal error; disconnect cancels generation and closes its stream.

\paragraph{Citation Engine.}
Positive integer labels are allocated within each request and map to exact paper,
chunk or SQL evidence snapshots. Screening preserves surviving labels; SQL evidence
appends without renumbering them. Before semantic validation, the engine checks label
syntax, membership and trailing placement. Deterministic rendering attaches corpus
and chunk IDs, author IDs, title, year, section, source text and a SHA-256 source
fingerprint; SQL citations resolve by query ID to the emitted table and its scope.
The UI citation drawer exposes this provenance. Syntactic validity alone does not
establish support, so evaluation keeps structural checks separate from human citation
and numerical-claim judgments. Pipeline timing fields include retrieval and guardrails;
provider token throughput and engine-only TTFT are measured separately by the benchmark.

\paragraph{Local WP2 verification.}
A local diagnostic on the recorded dataset revision tested the citation path for the
paper with corpus ID 215737187. Direct S2 retrieval returned 20 citing papers; all
20 IDs and publication years matched a separately exported bounded database reference.
All 80 retrieved passages belonged to those 20 citing papers. This verifies one
retrieval case; it does not measure answer quality or baseline-use accuracy.

A plain S3 author-ranking request completed in approximately 99.8 seconds in the local
setup. Its three returned author IDs and counts (32379: 7, 41141: 6, and 12839: 5)
matched an independently calculated reference for the same three seed papers and
dataset revision. One sentence was checked and retained. These counts describe the
recorded subset of the dataset rather than global citation totals. In a separate S2
application attempt, retrieval returned evidence but the request timed out before an
answer was emitted; a corpus-ID phrasing was also incorrectly classified as a write
request. These failures are not counted as completed answers.

Automated retrieval, sentence-streaming and evaluation checks passed 193 tests, with
three skipped, in one local run. A later focused S2 and metrics run passed 109 tests,
with one skipped. A separate 20-passage development sample had 17 assistant-reviewed
labels and three unresolved passages. The baseline matcher was refined using that
sample; results on it must not be reported as independent human-reviewed accuracy.
Reviewed citation precision and numerical exact-match results across scenarios remain
to be measured.
